%% ma: L10-B13-0006
%% lop: 10
%% chuong: 5
%% bai: 13
%% dang: 10.13.4
%% loai: TL
%% muc: TH
%% dap_an: "-"
%% nguon: "chuyen-de-cac-so-dac-trung-cua-mau-so-lieu-khong-ghep-nhom-toan-10-chuong-trinh-moi.pdf (GV. Phan Nhật Linh) (P02), Bài 02 Dạng 1 Bài tập tự luận 6"
%% kien_thuc: "Lớp 10 Bài 13 (trung vị)"
%% trang_thai: chinh-thuc
%% ghi_chu: "Tự duyệt theo yêu cầu giáo viên 22/10/2026."
\begin{ex}
Điểm thi môn Toán của lớp 10B được cho trong bảng sau:
\begin{center}
\begin{tabular}{|c|c|c|c|c|}\hline
Điểm thi & $6$ & $7$ & $8$ & $9$\\\hline
Tần số & $9$ & $18$ & $12$ & $6$\\\hline
\end{tabular}
\end{center}
Tìm trung vị của mẫu số liệu.
\loigiai{$n=9+18+12+6=45$ (lẻ) nên trung vị là giá trị thứ $23$ sau khi sắp xếp.

Tần số tích lũy: $9$; $27$; … nên giá trị thứ $23$ là $7$. Vậy $M_e=7$.}
\end{ex}
